\subsection{The support intersection map with projective resolution strata}

\begin{statement}{C.2.MH.1}{Proved}{The analytic support comparison used here}
Let $U$ be a smooth complex threefold, let $S\subset U$ be a
compact divisor, and suppose that an embedded resolution
$r:U'\to U$ is a sequence of blowups with smooth compact
centres in $S$, with $D=(r^{-1}S)_{\rm red}$ a simple
normal crossings divisor whose components are projective.
Let $V$ be a finite-monodromy rational local system on $U$.
Suppose that the cohomology of the punctured sufficiently
small fibre tubes around $S$, with coefficients in $V$,
vanishes in every degree. Then the mixed Hodge structure
on $H^k(S,V)$ obtained from a projective hyperresolution
is pure of weight $k$.

In this assertion the support-to-restriction morphism is
the actual morphism
\[
 R\Gamma_D(U',r^*V)\longrightarrow
 \underset{D\subset W}{\mathop{\rm colim}}\,
 R\Gamma(W,r^*V)=R\Gamma(D,r^*V).
 \tag{C.MH.1}
\]
Neither a K\"ahler metric on $U$ nor a projective
compactification of $U$ is needed for this comparison.
\end{statement}
\begin{proof}
First take constant rational coefficients. Write $D_I$
for the intersection of the components indexed by a
nonempty set $I$, and $D^{[p]}=\coprod_{|I|=p+1}D_I$.
All these spaces are smooth and projective. The ordinary
\v Cech de Rham complex on $D^{[\bullet]}$, with its
\v Cech weight filtration and its differential-form
filtration, is therefore a cohomological mixed Hodge
complex $K(D)$ for $\mathbb Q_D$. Its cohomology in
degree $k$ has weights at most $k$.

Here is the precise analytic construction of the other
complex and of the arrow. For each stratum $C=D_I$, put
\[
 K(C\log D)_{\mathbb C}
   =\Omega_{U'}^\bullet(\log D)\otimes\OO_C.
 \tag{C.MH.2}
\]
The tensor product in (C.MH.2) is a complex: in local
normal crossings coordinates, the ideal of $C$ is
preserved by the logarithmic differential. Endow it with
the quotient residue filtration and the filtration by
form degree. The quotient map from the restriction of
$\Omega_{U'}^\bullet(\log D)$ is a filtered
quasi-isomorphism to (C.MH.2). Its rational realization is
$i_C^*Rj_*\mathbb Q_{U'\setminus D}$, where $j$ is the
open inclusion. To see the comparison, apply the
iterated Poincar\'e residue: its residue-degree-$a$
quotient is a direct sum of the de Rham complexes of
$D_J\cap C$, shifted by $[-a]$ and with Tate twist
$(-a)$. The ideal and its differentials restrict these
forms to precisely $D_J\cap C$. This is the local
coordinate proof of \cite[Section~3, Theorem~1 and
Corollary~2]{Steenbrink1999}; it uses holomorphic
coordinates and the projectivity, or compact K\"ahler
property, of the displayed strata. It does not use a
compactification of their ambient neighbourhood.

Totalizing (C.MH.2) over $D^{[\bullet]}$ gives the
cohomological mixed Hodge complex $K(D\log D)$, with
rational realization $i_D^*Rj_*\mathbb Q_{U'\setminus D}$.
The natural inclusion $K(D)\to K(D\log D)$ is a
morphism of mixed Hodge complexes. Its mixed cone,
shifted by $[-1]$, computes cohomology with support in
$D$; the cone projection is exactly (C.MH.1). This is
the excision comparison of
\cite[Section~4, equation~(3)]{Steenbrink1999}, applied
on the compact divisor. All complexes just used are
defined on that divisor, and all purity checks are on
its projective strata.

For completeness, the weight bound on the source can
be checked without invoking an invariant-cycle theorem.
The closed-support Mayer--Vietoris complex, together
with the holomorphic Thom isomorphisms, gives
\[
 E_1^{-p,q}=
 \bigoplus_{|I|=p+1}
 H^{q-2p-2}(D_I,\mathbb Q)(-p-1)
 \ \Longrightarrow\ H_D^{q-p}(U',\mathbb Q),
 \qquad p\ge0.
 \tag{C.MH.3}
\]
The first differential is the alternating sum of Gysin
maps for $D_I\hookrightarrow D_{I\setminus\{i\}}$.
The comparison with the preceding mixed cone can be
verified by induction on the number of components.
For one smooth divisor the residue exact sequence
identifies
\[
 \operatorname{Cone}\bigl[\Omega_{U'}^\bullet
       \longrightarrow\Omega_{U'}^\bullet(\log D_i)\bigr][-1]
       \simeq\Omega_{D_i}^\bullet[-2](-1),
\]
where $(-1)$ records both the shift of the Hodge
filtration by one and the rational $2\pi i$ residue.
For a union $A\cup B$, use the actual support triangle
$R\Gamma_{A\cap B}\to R\Gamma_A\oplus R\Gamma_B
\to R\Gamma_{A\cup B}$. Iterated residues commute
with its inclusion maps, with the alternating signs
fixed by the ordering of the omitted coordinates.
Induction identifies its totalization with the same
mixed cone. These comparisons agree on overlaps,
since they are the residue maps for the holomorphic
divisor line bundles. Thus (C.MH.3) compares the actual support
complex, including its arrow to $K(D)$, rather than
only its unfiltered vector spaces. Its displayed term
has pure weight $q$. In total degree $k=q-p$ this is
weight $k+p\ge k$. It follows that $H_D^k(U',\mathbb Q)$
has weights at least $k$.

Finite coefficients are handled by taking a finite
\'{e}tale cover on which $V$ is constant and then a
rational direct summand of a finite direct sum of its
regular representation.
The inverse images of the $D_I$ are still projective:
a finite holomorphic \'{e}tale cover of a
projective manifold is projective, since the pullback
of an ample bundle is positive. The deck action acts
on the displayed mixed Hodge complexes. Rational
idempotents consequently give exactly the same
comparison and weight bounds with coefficients in $V$.

By the assumed punctured-tube acyclicity, the support
long exact sequence makes (C.MH.1) an isomorphism on
cohomology. Properness of $r$ identifies the punctured
tubes upstairs and downstairs, since its centres lie
in $S$. The map is a morphism of mixed Hodge structures,
whose source has weights at least $k$ and whose target
has weights at most $k$. Strictness therefore proves
that $H^k(D,r^*V)$ is pure of weight $k$.

We finally make the comparison with the original
support explicit. The proper degree-one map between
the smooth ambient manifolds has the trace map, defined
by oriented Poincar\'e duality and proper pushforward,
and the projection formula gives
\[
 r_*r^*=\operatorname{id}:H^k(U,V)\longrightarrow H^k(U,V).
 \tag{C.MH.4}
\]
The same identity holds on every smaller ambient
neighbourhood. Passing to the direct limit of these
identities gives an injection
\[
                  r^*:H^k(S,V)\hookrightarrow H^k(D,r^*V).
 \tag{C.MH.5}
\]
The identifications in this limit are the usual
continuity identifications for a compact triangulable
subset. They can equally be made using cofinal
triangulated regular neighbourhoods retracting onto
$S$ and $D$. Properness ensures that inverse images of
neighbourhoods are cofinal; thus the map in (C.MH.5)
is the pullback by $D\to S$.

This last map is a mixed Hodge morphism as well. One
may construct the projective hyperresolution of $S$
from its projective normalization resolution and
resolutions of the conductor and exceptional curves.
To compare it with $D^{[\bullet]}$, resolve their
fibre products over $S$. Each such fibre product is a
closed analytic subspace of a product of projective
spaces already occurring in the two resolutions;
Chow's theorem makes it algebraic and projective,
and its resolution is projective. The resulting
common hyperresolution realizes (C.MH.5) by pullbacks
between projective strata. Cohomological descent then
identifies this with the stated mixed Hodge structures.
Injectivity into the pure target proves the assertion.
\end{proof}



\begin{statement}{C.2.MH.2}{Proved}{Euler number from a ruled normalization}
Let $f:Y\to\PP^1$ satisfy the full CDP hypotheses
$b_1(Y)=b_2(Y)=0$ with smooth torus general fibre.
Let $mS$ be a full fibre with irreducible reduced support,
and let $Z\to T\to S$ be a resolution of its
normalization. If $Z$ is rational, then
\[
                              e(S)=2.
 \tag{C.MH.6}
\]
If $Z$ is ruled over a curve of genus $g\ge1$, then
\[
                              e(S)=2-2g.
 \tag{C.MH.7}
\]
The calculation retains all multiplicities in the full
fibre and places no normal crossings assumption on $S$.
\end{statement}
\begin{proof}
The integral class of D.1.9 restricts nontrivially to
$S$ and, by D.1.8, nontrivially to the nearby torus.
More precisely, restrict the section
$\gamma\in H^0(\PP^1,R^1f_*\mathbb Z)$ to a small
disc $\Delta$ about this fibre. The integral Leray edge
isomorphism
$H^1(f^{-1}\Delta,\mathbb Z)
 \simeq H^0(\Delta,R^1f_*\mathbb Z)$ supplies the
actual integral class on the tube: the two adjacent
base groups $H^1(\Delta,\mathbb Z)$ and
$H^2(\Delta,\mathbb Z)$ vanish.
Choose a prime $\ell$ for which its restriction to that
torus is nonzero modulo $\ell$. On a sufficiently small
fibre tube $U$, its reduction defines a connected cyclic
\'{e}tale cover with group $\mathbb Z/\ell$. Let $V$ be
the rational augmentation local system, that is, the
kernel of the trace from the regular representation
to the constant representation. Its rank is
$d=\ell-1$, it is self-dual, and $H^0(S,V)=0$.
Over $\mathbb C$ it is the sum of the nontrivial
characters of this cyclic group. Since $\ell$ is prime,
each of these characters restricts nontrivially to
the nearby torus $F$. The group-cochain complex for
$\mathbb Z^4$ with such a character is acyclic: one of
the four commuting operators $\rho(e_i)-1$ is
invertible, and its inverse gives a contraction of the
Koszul complex. Thus
\[
 H^j(F,V)=0\quad\hbox{for every }j,
 \qquad H^j(U\setminus S,V)=0\quad\hbox{for every }j.
 \tag{C.MH.8}
\]
The second assertion is the Leray sequence of the
punctured-disc torus bundle. This argument uses the
smooth fibres of the original map, so it also applies
when the full central fibre is $mS$.

Choose an embedded resolution of $(U,S)$. Its strict
transform is a blowup of $Z$, hence projective. Every
other exceptional surface is obtained by blowups from
a projective bundle over a compact smooth curve, or
from $\PP^2$. It is therefore projective. All centres
are compact and lie in the central fibre. C.2.MH.1
applies and gives
\[
                         H^3(S,V)\text{ pure of weight }3.
 \tag{C.MH.9}
\]

Let $A\subset T$ and $\Gamma\subset S$ be the reduced
conductor inverse image and image. The finite
normalization square is an actual topological pushout.
Its cohomological descent sequence, with the pulled
back local system, gives an exact sequence of mixed
Hodge structures
\[
\begin{gathered}
0\longrightarrow C_V\longrightarrow H^3(S,V)
  \longrightarrow H^3(T,V)\longrightarrow0,\\
C_V=\operatorname{coker}\bigl[H^2(T,V)\oplus H^2(\Gamma,V)
                         \longrightarrow H^2(A,V)\bigr].
\end{gathered}
 \tag{C.MH.10}
\]
This also covers the normal case, with the conductor
terms zero. The left term is pure of weight two:
$H^2(A,V)$ is the top cohomology of a compact curve
with a finite local system, and hence pure of weight
two. Strictness and (C.MH.9) make that term zero.
Consequently
\[
                           H^3(S,V)\simeq H^3(T,V).
 \tag{C.MH.11}
\]

For the resolution $\pi:Z\to T$, the restriction
$H^2(Z,V)\to H^2(E,V)$ to its exceptional curve is
surjective. Indeed $V$ is constant on each exceptional
fibre, since it is pulled back from $T$; the Gysin
classes of its irreducible exceptional curves restrict
by their negative-definite intersection matrix. That
matrix is invertible over $\mathbb Q$, separately
for each exceptional fibre and each coefficient
vector. The resolution descent sequence therefore
gives
\[
 H^3(T,V)\simeq H^3(Z,V),\qquad
 H^4(S,V)\simeq H^4(T,V)\simeq H^4(Z,V).
 \tag{C.MH.12}
\]
The last two identifications also follow directly
because all the omitted loci have real dimension two.

There is a second, purely topological consequence of
(C.MH.8). The support sequence and oriented local
Poincar\'e duality give
\[
 H^k(S,V)\simeq H_S^k(U,V)
             \simeq H_{6-k}(S,V).
 \tag{C.MH.13}
\]
Again these are the natural direct-limit maps on
sufficiently small tubes. Since $S$ has real
dimension four, (C.MH.13) gives $H^1(S,V)=0$ and,
using the self-duality of $V$,
$h^2(S,V)=h^4(S,V)$. Therefore (C.MH.11)--(C.MH.13)
give the exact numerical identity
\[
 d\,e(S)=\chi(S,V)=2h^4(Z,V)-h^3(Z,V).
 \tag{C.MH.14}
\]

If $Z$ is rational, the pullback of $V$ is constant,
$h^4(Z,V)=d$, and $h^3(Z,V)=0$; this proves (C.MH.6).
If $Z$ is ruled over $B$, blowups leave its fundamental
group and its third cohomology unchanged, and the
pulled-back local system comes from $B$. The
projective-bundle formula gives
$H^3(Z,V)=H^1(B,V)(-1)$. If the cyclic character is
trivial on $B$, its dimensions are
$h^4(Z,V)=d$ and $h^3(Z,V)=2gd$. Otherwise every
nontrivial character of the prime-order group remains
nontrivial on $B$, so $h^4(Z,V)=0$ and
$h^3(Z,V)=(2g-2)d$. Both substitutions in (C.MH.14)
give (C.MH.7).
\end{proof}

\begin{statement}{C.2.MH.3}{Proved}{The contracted irrational section contributes zero}
In the contracted-horizontal-section situation of
C.2.CR.1, the normalization resolution is ruled over
an elliptic curve and
\[
                    e(S)=0,\qquad n=v-N_0.
 \tag{C.MH.15}
\]
\end{statement}
\begin{proof}
C.2.CR.5 proves that the ruling genus is one.
C.2.MH.2 then gives $e(S)=0$ for the actual fibre.
The exact equality (C.CR.3) gives $n=v-N_0$.
The later unconditional calculation C.2.UN.7
gives $e(S)\ge2$ and $\pi_1(S)=1$ for this
contracted irrational configuration. Thus its
comparison with the necessary equations here
excludes the case entirely under the CDP
hypotheses; this reduction supplies no example
of an actual CDP fibre.
\end{proof}

\begin{statement}{C.2.MH.4}{Proved}{Every irreducible class VII normalization contributes zero}
Under the full CDP torus-fibration hypotheses, an
irreducible fibre support whose normalization resolution
is a surface of class VII has
\[
                              e(S)=0.
 \tag{C.MH.16}
\]
Consequently every such actual fibre satisfies the
equality data of C.2.VII.4.
\end{statement}
\begin{proof}
Use the prime-order rational augmentation local system
$V$ of C.2.MH.2, of rank $d=\ell-1$. Its construction,
the torus acyclicity (C.MH.8), and the support duality
(C.MH.13) do not require a projective normalization.
They give $h^0(S,V)=h^1(S,V)=0$ and
$h^2(S,V)=h^4(S,V)$. Conductor descent still gives the
surjection $H^3(S,V)\to H^3(T,V)$ and the isomorphism
$H^4(S,V)\simeq H^4(T,V)$. The negative-definite
exceptional Gysin matrix in the proof of (C.MH.12)
is valid for the resolution of any normal complex
surface. It gives
\[
 \begin{gathered}
 h^3(S,V)\ge h^3(Z,V),\qquad h^4(S,V)=h^4(Z,V),\\
 d\,e(S)\le 2h^4(Z,V)-h^3(Z,V).
 \end{gathered}
 \tag{C.MH.17}
\]
No mixed Hodge assertion is used in these comparisons.

If the cyclic character restricts nontrivially to $Z$,
all its nontrivial powers remain nontrivial because
$\ell$ is prime. Poincar\'e duality on the compact
oriented surface gives $h^4(Z,V)=0$, and (C.MH.17)
gives $e(S)\le0$. If its restriction is trivial,
$h^4(Z,V)=d$ and $h^3(Z,V)=d$: a class VII surface
has $b_1=b_3=1$, and point blowups preserve $b_3$.
Thus $e(S)\le1$ in this second case. C.2.UN.1 proves
that the Euler number of every actual torus-degeneration
fibre is even, so the second case also gives
$e(S)\le0$. The nonnegative lower bound of C.2.VII.3
now gives (C.MH.16) in both cases. C.2.VII.4 supplies
its exact normalization-and-conductor equality data.
\end{proof}
